% Generated draft: requires mathematical and source review.
\documentclass[11pt]{article}
\usepackage[margin=25mm]{geometry}
\usepackage{fontspec}
\setmainfont{Latin Modern Roman}
\usepackage{amsmath,amssymb,mathtools,unicode-math}
\setmathfont{Latin Modern Math}
\usepackage{enumitem,needspace,xurl,hyperref}
\setlength{\parindent}{0pt}
\setlength{\parskip}{5pt}
\setlength{\emergencystretch}{3em}
\newcommand{\N}{\mathbb{N}}
\newcommand{\Z}{\mathbb{Z}}
\newcommand{\Q}{\mathbb{Q}}
\newcommand{\R}{\mathbb{R}}
\newcommand{\C}{\mathbb{C}}
\newcommand{\F}{\mathbb{F}}
\newcommand{\A}{\mathbb{A}}
\newcommand{\PP}{\mathbb{P}}
\newcommand{\E}{\mathbb{E}}
\newcommand{\eps}{\varepsilon}
\newcommand{\dd}{\,\mathrm{d}}
\begin{document}
\section*{UnsolvedMath ID: 20003271}
\textbf{Source problem code:} catalogue entry
\subsection{Definitions and mathematical statement}
The printed Conjecture 5 concerns a connected semisimple Lie group $G$ whose every simple factor has real rank at least two, or a lattice in such a group. It claims that every action having at least one nontrivially uniformly partially hyperbolic element is $C^\infty$-conjugate to an algebraic action. The source itself has the grammatical errors ``which that'' and ``partially hyperbolicly''; the supplied verification labels their repair as editorial, not OCR. The final fragment ``Partial results on'' has no recoverable continuation.

A precise formulation still needs the ambient manifold, compactness, smoothness of the action, the meaning of nontrivial partial hyperbolicity and the class of algebraic models. In the aid's compact smooth setting, uniform partial hyperbolicity uses an invariant dominated splitting $TM=E^s\oplus E^c\oplus E^u$ with uniform stable contraction and unstable expansion and nonzero hyperbolic directions. The conjecture asks for one smooth conjugacy simultaneously intertwining every acting group element, not merely a conjugacy of the distinguished element.

The aid supplies a factor-hiding construction. Let $\Gamma_2$ be any higher-rank lattice and $\beta:\Gamma_2\to\operatorname{Diff}^\infty(N)$ any smooth action on a compact manifold $N$. Then
\[
 \alpha:\bigl(SL(3,\Z)\times\Gamma_2\bigr)\curvearrowright\mathbb T^3\times N,
 \qquad
 \alpha(\gamma_1,\gamma_2)(x,y)=(\gamma_1x,\beta(\gamma_2)y)
\]
is a smooth action of a product higher-rank lattice. Choose a hyperbolic $A\in SL(3,\Z)$. The element $(A,e)$ acts as $A\times\mathrm{id}_N$, with the toral stable and unstable subbundles and center $TN$, and is uniformly partially hyperbolic. The aid mentions an explicit matrix but supplies none, so the statement uses the stated hyperbolic-matrix hypothesis rather than fabricating its example.

No dynamical restriction on $\beta$ is needed for that one-element premise. If $\beta$ is faithful, the product action is faithful. If it preserves a volume, the product preserves Haar volume times that volume; if the $\Gamma_2$ action is ergodic for it, the full product action is reported ergodic. These are conditional preservation statements, not assertions that every arbitrary $\beta$ has these properties. For a nontrivial connected positive-dimensional $N$, stable/unstable paths for the distinguished element stay in $\mathbb T^3\times\{y\}$, so its accessibility fails across the $N$ direction. Full-action ergodicity therefore does not imply accessibility of that element.

A parallel construction is reported for connected product groups, using a regular right-translation action on a compact quotient for the hyperbolic factor and an arbitrary smooth action of the other factor. Exact quotient and regular-element hypotheses are not supplied. Both constructions show that one partially hyperbolic element can be confined to one factor while another factor acts arbitrarily; they do not establish a definition-independent obstruction to algebraicity. The aid explicitly declines to call this a counterexample because algebraic and generalized quasi-affine normal forms, including permissible auxiliary fibers, are not fixed. The broad conjecture remains open in the supplied account, rather than disproved merely by the product formula.
\subsection{Short English statement}
The aid shows why one partially hyperbolic element is a weak premise for reducible groups: it can act only on one product factor while another factor acts arbitrarily. Faithfulness, volume preservation and full-action ergodicity do not automatically remove that issue. No precise nonalgebraicity invariant or counterexample is supplied.
\subsection{Sources}
\begin{itemize}
\item Local UnsolvedMath registry, permanent ID 20003271, source code catalogue entry; \texttt{aim-other-notes.json}, notes[77], source tag conjecture.
\item AIM, \emph{Emerging applications of measure rigidity}, Conjecture 5: \url{https://aimath.org/WWN/measrigid/measrigid.pdf}. The supplied verification says both grammatical errors occur in the PDF; this draft does not independently repeat the check.
\item The factor-hiding lattice and connected-group constructions appear in the supplied research summary without separate full bibliographic references or its advertised explicit matrix.
\item Machine-generated research attempt 1: reduction, partially solved at the run, rigor labelled ``proved,'' low novelty confidence. Not peer reviewed; independently verify before citation. Registry updated 14 August 2026.
\end{itemize}
\end{document}
